\section{Open Questions}

Five genuine open questions arose from re-reading the paper and the audit
report. None are closable without new data or new modelling; all are
recommended follow-ups for the LUCID-100 programme.

\begin{enumerate}
  \item \textbf{Neutron-energy dependence of DSB yield.} Is the DSB-yield-per-Gy
        for p(66)/Be(40) fast neutrons in human lymphocytes energy-dependent
        within the neutron spectrum? The paper does not include the spectrum
        (nominal mean $\sim$29~MeV) and does not compare against 14~MeV D-T or
        62~MeV clinical proton-neutron secondary fields.
        \emph{Next steps:} extract iThemba spectrum from dosimetry refs;
        cross-check against ICRU 63/95 reference fields; compare per-Gy DSB
        yield to published 14~MeV (JAERI, PTB) and 62~MeV (Louvain, HIT)
        lymphocyte $\gamma$-H2AX datasets.

  \item \textbf{RBE vs LET curve reproducibility across labs.} The paper does not
        calculate an internal RBE; any RBE inference relies on cross-study
        low-LET baselines that may differ in donor pool, scoring method (manual
        vs Metafer), or fixation time.
        \emph{Next steps:} assemble published low-LET ($^{60}$Co, 6~MV, 250~kVp)
        $\gamma$-H2AX per-cell foci-per-Gy in lymphocytes at 30~min post-IR;
        compute RBE$_{\text{DSB}}$ with covariate correction; compare against
        MCDS/PARTRAC predictions given the p(66)/Be(40) spectrum.

  \item \textbf{Translation to modern PBL biodosimetry protocols.} Does the
        $\sim$40\,\% HDR-over-LDR DSB-yield effect translate to the actual
        clinical biodosimetry endpoints (dicentric assay, PCC, micronuclei),
        or is it a $\gamma$-H2AX-specific phenomenon (e.g. attributable to
        foci merger / scoring saturation at high induction rates)?
        \emph{Next steps:} meta-analyse published dicentric/micronucleus
        yields for fast neutrons at HDR vs LDR in lymphocytes; if a per-cell
        foci CSV becomes available, estimate foci-merger correction; test
        whether the K-coefficient peak at 0.250~Gy reflects a
        scoring-saturation artifact.

  \item \textbf{Mixed field (neutron + gamma) DSB response.} Real high-LET
        exposures (BNCT, reactor, criticality) always include a substantial
        gamma component. Whether the mixed-field DSB yield is predictable
        linearly from pure-neutron + pure-gamma responses, or shows
        super-additive damage/repair coupling, is unresolved by this paper.
        \emph{Next steps:} search BNCT lymphocyte $\gamma$-H2AX studies
        (Kyoto, Petten, Argentina) with neutron-fraction-resolved DSB yields;
        compute (n+$\gamma$) $-$ predicted-additive on any matched dataset;
        model with MCDS coupled neutron and gamma track structure.

  \item \textbf{Coupling to cytogenetic aberration models.} The paper stops at
        DSB yield and repair kinetics; it does not link these to chromosomal
        aberration formation, which is the endpoint that matters for
        biodosimetry and risk. Can the fitted induction (poly2) and repair
        (single-exp) parameters be plugged into a two-lesion kinetic (TLK) or
        Sachs-Rossi repair-misrepair model to predict dicentric yield at HDR
        vs LDR?
        \emph{Next steps:} parameterise a TLK model with the paper's fitted
        rates; predict per-dose dicentric yield; compare against Vral / IAEA
        fast-neutron dicentric-vs-dose curves.
\end{enumerate}
